% !PRINTED 161
\weberpage{161}
% [source-original mathematical slip, confirmed 2026-08-25 against scan p. 186 / printed p. 161]
% The print assigns degrees m and n, respectively, to f and phi, although printed p. 156
% defined f with degree n and phi with degree m. The facing translation preserves that order.
We can then also regard $f$ and $\varphi$ as homogeneous functions
of degree $m$ and $n$ in the variables $x, y_1, z_1$, and
the resultant then becomes a homogeneous function $R(y_1, z_1)$
of degree $mn$ in $y_1$ and $z_1$. If we set in it

\[
y_1 = b' + b\lambda, \qquad z_1 = c' + c\lambda,
\]

we obtain an equation of degree $mn$ in $\lambda$, and,
provided $R(y_1, z_1)$ does not vanish identically, we can choose the constants
$b, b', c, c'$ so that the coefficient of the highest
power of $\lambda$, namely $R(b, c)$, is different from zero. Then
to each root of the resultant there corresponds only one system of values
of the ratio $x:y:z$, and so the number of common
roots of $f$ and $\varphi$ is at most equal to $mn$. In order to be able to
state the theorem that the number of common roots
is always equal to $mn$, one must further adopt a convention
according to which certain of these pairs of values are to be counted multiply.

The theorem which, in geometric dress, says that
two curves of the $m$th and $n$th order intersect in $mn$ points
is called \gesperrt{Bezout's theorem}.

\begin{center}
§.\,50.

\gesperrt{Elimination from several equations.}
\end{center}

The principles of elimination, which we have applied in the
foregoing to two equations, can be extended without
difficulty to several equations with a corresponding
number of unknowns.

Let us first take the equations

\weq{1}{f(x) = 0, \quad \varphi(x) = 0, \quad \psi(x) = 0,}

of degrees $m, n, p$, whose coefficients $a_i, b_i, c_i$ are at first to be
regarded as independent variables.

We apply to two of them, say to $f(x), \varphi(x)$, the
algorithm of the greatest common divisor and
obtain a last remainder, which is a rational function of the $a$
and $b$ and which, freed of denominators, agrees with the resultant $R$
of $f(x)$ and $\varphi(x)$, and a last remainder but one, which is a
\gesperrt{linear function} of $x$. The condition for the existence
of a common root of $f = 0$ and $\varphi = 0$
is the vanishing of the resultant, and the last remainder but one
%CONT>>
